% Mathematical reductions for the implementation introduced after v0.5.2.
\section{Reducing the inference problem}
\label{sec:reduction}

Let $q=\h$. The efficient calculation begins by distinguishing information
that is absent, information that is exact, and information that remains an
interval. These distinctions concern the observation model, not a different
estimand. Throughout this section the censoring mixture is excluded. Its
sequential interpretation remains Algorithm~M.

\subsection{Marginalize absent observations; condition exact ones}

Partition a jointly Gaussian vector into unobserved non-target coordinates
$U$, pinned coordinates $P$ observed at $p$, and retained coordinates $X$.
The latter contain the requested target and the remaining intervals.
Integrating out $U$ simply selects the $(X,P)$ principal covariance block;
it is not conditioning $U$ to zero. For a zero-mean prior, exact conditioning
then gives
\begin{align}
 m_X &= \Sigma_{XP}\Sigma_{PP}^{-1}p,\label{eq:pin-mean}\\
 V_X &= \Sigma_{XX}-\Sigma_{XP}\Sigma_{PP}^{-1}\Sigma_{PX}.
 \label{eq:pin-var}
\end{align}
The inverse notation denotes linear solves. A singular pin block requires
compatible observations in its Gaussian support, not arbitrary jitter.
A zero conditional variance denotes a deterministic coordinate: its bound
must contain that value. Incompatible exact observations are errors.

The covariance in \eqref{eq:pin-var} depends on the pedigree covariance and
pin mask, not the numerical onset values. Families with the same masks can
share the reduction. With no remaining interval, Gaussian conditioning is
the whole calculation. With one remaining interval, the univariate selection
formula is exact. With several intervals, apply PA to the centered remaining
bounds, preserving their relative order. Conditioning pins first changes the
old PA approximation when a pin previously followed an interval. It leaves
the specified joint distribution unchanged; it does not prove that every
multi-interval approximation becomes more accurate.

For family-free ADuLT, let $m,v$ be the scalar mean and variance of the full
liability conditional on its observation. Gaussian regression gives
\begin{equation}
 \E[a_o\mid D_o]=qm,\qquad
 \Var(a_o\mid D_o)=q(1-q)+q^2v.
 \label{eq:adult-exact}
\end{equation}
A pinned value $t$ gives $qt$ and $q(1-q)$. No family covariance or numerical
integration is necessary. With fixed bounds and $q>0$, changing $q$ only
rescales this family-free phenotype. An ordinary scale-invariant linear or
mixed-model association statistic therefore does not require estimating $q$
for this use. Family-informed scores and liability-scale reporting do.

\subsection{Nuclear families have at most two parental factors}
\label{sec:parent-factors}

Assume unrelated, noninbred parents, the additive model, and independent
residual environments. Write $z=(a_m,a_f)^\mathsf{T}\sim N(0,qI_2)$.
For offspring $j$, including the proband,
\begin{equation}
 a_j=\tfrac12(a_m+a_f)+M_j,\qquad
 M_j\overset{\mathrm{ind}}\sim N(0,q/2).
 \label{eq:offspring-factor}
\end{equation}
The Mendelian terms are independent of $z$. Conditional on $z$, all observed
full liabilities are independent normals,
\begin{equation}
 \ell_j\mid z\sim N(b_j^\mathsf{T}z,d_j),\qquad
 (b_m,b_f,b_{\rm child})=((1,0)^\mathsf{T},(0,1)^\mathsf{T},
                          (1/2,1/2)^\mathsf{T}),
 \label{eq:factor-liability}
\end{equation}
with $d_m=d_f=1-q$ and $d_{\rm child}=1-q/2$.
A missing observation contributes one. An interval contributes
\begin{equation}
 L_j(z)=\Phi\!\left(\frac{u_j-b_j^\mathsf{T}z}{\sqrt{d_j}}\right)
       -\Phi\!\left(\frac{l_j-b_j^\mathsf{T}z}{\sqrt{d_j}}\right),
 \label{eq:factor-likelihood}
\end{equation}
whereas a pin contributes the corresponding normal density at its observed
value. Thus
\begin{equation}
 p(z\mid D)\ \propto\ \phi_2(z;0,qI_2)\prod_j L_j(z).
 \label{eq:factor-posterior}
\end{equation}
The number of siblings changes the number of scalar factors, not the
integration dimension. If no parental observations remain, their shared mean
$S=(a_m+a_f)/2\sim N(0,q/2)$ suffices and the integral is one-dimensional.
Personalized thresholds do not alter this factorization.

Pins can be absorbed before integration. For $q>0$ and positive $d_j$, the
conditional factor precision and mean are
\begin{equation}
 H=(qI_2)^{-1}+\sum_{j\in P}\frac{b_jb_j^\mathsf{T}}{d_j},\qquad
 m_z=H^{-1}\sum_{j\in P}\frac{b_jp_j}{d_j}.
 \label{eq:factor-pins}
\end{equation}
Only unpinned interval factors remain. Integrating in a Gaussian coordinate
system adapted to the posterior mode further prevents rare observations
from placing most posterior mass beyond the original prior-centered nodes.
Products and normalizing constants are evaluated in log space.

To recover the proband, let $S=b_o^\mathsf{T}z$, $v_o=1-q/2$, and
$r=q/2$. Let $m_o(z),v_o(z)$ be the scalar moments of $\ell_o$ after its own
observation at fixed $z$; for an unobserved proband these are $S,v_o$.
Then
\begin{align}
 \mu_a(z)&=S+\frac{r}{v_o}\{m_o(z)-S\},\label{eq:factor-target-mean}\\
 V_a(z)&=r-\frac{r^2}{v_o}+\frac{r^2}{v_o^2}v_o(z).
 \label{eq:factor-target-var}
\end{align}
The requested mean is $\E_{z\mid D}[\mu_a(z)]$. Its posterior variance is
\begin{equation}
 \Var(a_o\mid D)=\E_{z\mid D}[V_a(z)]
                  +\Var_{z\mid D}[\mu_a(z)].
 \label{eq:factor-total-var}
\end{equation}
The first term must not be dropped. Full-liability moments follow by using
$m_o(z),v_o(z)$ instead. With $Q$ nodes per factor and $s$ observations,
the two-dimensional quadrature costs $O(sQ^2)$, with $O(Q^2)$ working storage
when observations are processed one at a time. Increasing the resolution
provides a convergence diagnostic for both moments; agreement of successive
resolutions is not a certified error bound. The implementation rejects
unresolved convergence rather than returning an apparently precise score.

This is a factorization of the additive nuclear-family model, not a new
observation model. The initial quadrature interface excludes shared
C/M environments, arbitrary/inbred pedigrees, multiple traits, and the
censoring mixture. It requires $q<1$ because a parental residual becomes
singular at $q=1$; $q=0$ is an analytic special case. PA and Gibbs retain
their broader model support.

\subsection{Compute selected relationships with bounded reuse}

Ancestral closure is needed to calculate exact relatedness, but its entire
dense relationship matrix is not needed for a score. With parents $s_i,d_i$
and a parent-before-child ordering, the tabular identities are
\begin{equation}
 A_{ii}=1+\tfrac12A_{s_i d_i},\qquad
 A_{ij}=\tfrac12(A_{s_i j}+A_{d_i j})\quad(i\text{ later than }j).
 \label{eq:selected-kinship}
\end{equation}
An unknown parent contributes zero; distinct unrelated founders have zero
relationship. Evaluate \eqref{eq:selected-kinship} only for requested pairs
and their dependencies, retaining the complete ancestry graph. An explicit
stack avoids recursion-depth limits. A bounded cache shares computed entries
between overlapping probands; eviction changes work, not the result.
This saves repeated local dense matrices but does not imply linear worst-case
time: the number of required pairs can still approach the square of the
population size. Transient dependency storage also depends on pedigree depth.

Preserve the current observation order when reducing a PA input. Preserve
the observation counts and calendar-time masking separately from inference
coordinates. Near singular parameter boundaries, covariance repair can
change the effective model; the driver retains the full construction when
its conservative positive-definiteness bound cannot rule out that repair.

\subsection{Fit covariance parameters by pairwise likelihood}
\label{sec:pairwise-theory}

For common binary thresholds and independent families, parameter estimation
need not reconstruct every latent liability. Let $\theta=(q,c^2,m^2)$,
restricted to identifiable nonnegative components with positive residual
variance. For a pair,
\begin{equation}
 \rho_{ij}(\theta)=qA_{ij}+c^2C_{ij}+m^2M_{ij}.
 \label{eq:pairwise-rho}
\end{equation}
Let $p_{ab}(t,\rho)$ be its bivariate-normal probability for binary outcomes
$a,b\in\{0,1\}$ at threshold $t$. The composite log likelihood is
\begin{equation}
 \ell_c(\theta)=\sum_f w_f\sum_{(i,j)\in f}
           \log p_{y_i y_j}\{t,\rho_{ij}(\theta)\}.
 \label{eq:composite-likelihood}
\end{equation}
The weight is one for population sampling; valid positive inverse inclusion
probabilities can supply $w_f$ under the existing IPW contract. Pairs sharing
the same threshold and component design have the same four probabilities,
so their weighted $2\times2$ tables can be accumulated once. This compresses
\emph{the composite objective}; it does not make pair counts sufficient for
the full multivariate family likelihood.

At a common threshold, a stable expression for a discordant cell is
\begin{equation}
 p_{01}(t,\rho)=\frac{1}{2\pi}\int_{\arcsin\rho}^{\pi/2}
       \exp\!\left\{-\frac{t^2}{1+\sin u}\right\}\,du.
 \label{eq:pair-discordance}
\end{equation}
The other cells follow from the known marginal probabilities. This avoids
subtracting two nearly equal bivariate probabilities as $\rho$ approaches
one. Deterministic quadrature and analytic probability derivatives permit
low-dimensional constrained optimization \cite{renard2004,bravo2019}.

Within-family pairs are dependent. At an identifiable interior solution,
let $s_f$ be the family score, $H=-\nabla^2\ell_c$, and $J$ the empirical
outer-product sum of weighted family scores. The sandwich covariance is
\begin{equation}
 \widehat{\Var}(\widehat\theta)=H^{-1}JH^{-1}.
 \label{eq:pairwise-sandwich}
\end{equation}
The implementation centers family scores and applies the finite-cluster
correction. It does not report the ordinary inverse Hessian as the sampling
covariance. Boundary or singular-information fits do not justify an interior
normal approximation and return explicitly unavailable standard errors.
Family bootstrap is appropriate only under its independent-cluster sampling
assumptions. Ascertainment is not repaired by changing the numerical engine;
personalized/pinned thresholds and overlapping extracted register families
remain outside this initial fitter's contract.

Pairwise fitting estimates the same covariance parameters with a different
statistical criterion from the stochastic moment fitter. Neither matching
point estimates on a few simulations nor passing probability tests establishes
coverage or statistical efficiency. The new method is explicit and opt-in.

\subsection{Implementation pilot and its limits}
\label{sec:efficiency-pilot}

The v0.6.0 development pilot of 5 September 2026 preserves its measured
sources, input arrays and outputs under
\texttt{benchmarks/results/2026-09-05-efficient-inference/}.
The composite source hash begins \texttt{67002dd28513}; the base commit is
\texttt{8856ba8}, with the implementation changes included in the snapshot.
The machine was an Apple M2 Pro (10 cores, 16 GB), macOS 26.6.2;
Python 3.10.20, NumPy 2.2.6, SciPy 1.15.3 and Numba 0.67.0.
AC power and Low Power Mode off were checked before and after.
Numba used one thread; requested BLAS thread limits and build configuration
are recorded, but runtime BLAS thread counts were unavailable.

\begin{table}[ht]
\centering
\caption{Small implemented workloads, not a population-scale extrapolation.
Warm time is the median of three calls. Each arm runs in a separate process
with an empty initial Numba cache. First-call time includes compilation
reached by that call; neither call timing includes imports. Peak RSS covers
the whole process and all repetitions.}
\label{tab:efficient-pilot}
\small
\begin{tabular}{@{}lrrr@{}}
\toprule
Workload and method & First call (s) & Warm (ms) & Peak (MiB) \\
\midrule
One family: PA & 1.060 & 0.223 & 185.4 \\
One family: quadrature & 0.604 & 1.929 & 171.5 \\
One family: Gibbs, $10^6$ draws & 2.314 & 465.804 & 211.0 \\
\midrule
20,000 ADuLT: scalar & 0.721 & 0.398 & 204.3 \\
20,000 ADuLT: covariance API & 1.374 & 11.391 & 203.1 \\
\midrule
200 probands: selected relationships & 1.091 & 59.101 & 190.4 \\
200 probands: full construction & 1.317 & 250.229 & 186.2 \\
\bottomrule
\end{tabular}
\end{table}

The nuclear example uses $q=0.5$, $K=0.05$, own liability pinned at 2,
a control mother, affected father, one affected and three unaffected siblings.
Quadrature gives mean $1.2390361132$ and variance $0.2131471559$;
Gibbs gives $1.2388572910$ and $0.2131210024$, with mean Monte Carlo SE
$0.000147283$. Their means differ by $1.21$ Monte Carlo SEs.
Quadrature stops at 64 nodes per active dimension; its last-two-refinement
change is $4.44\times10^{-15}$. That diagnostic does not bound the true error.
Current pin-first PA gives $1.2452930026$ and $0.2150592363$ and remains
the fastest method here. Its remaining-interval approximation is visible
despite exact pin conditioning.

ADuLT uses equal groups of pinned, case, control and absent observations at
$q=0.5$. Scalar and covariance means are identical; the largest variance
difference is $1.11\times10^{-16}$. The register workload has 461 people,
a shared 60-generation paternal chain, degree 1 and 200 probands.
Both arms use the current PA engine: only relationship construction and
coordinate selection differ. Means and variances are identical.
The selected path is $4.23\times$ faster on warm calls, but peak RSS is
slightly higher; this is no evidence for lower memory use at population scale.

Three population-sampled cohorts of 1,000 nuclear families, $q=0.5$, $K=0.1$,
give pairwise estimates $(0.636,0.507,0.465)$ and moment estimates
$(0.755,0.412,0.482)$ on identical inputs. Pairwise calls take
0.057--0.058 seconds; the moment fitter's default 1500 iterations, 500 burn-in
and five inner sweeps take 2.379--2.421 seconds. The criteria differ, and
default controls do not establish convergence. These figures therefore do
not compare cost at equal accuracy, sampling efficiency or interval coverage.
